\section{Introduction}

According to the Sapir-Whorf hypothesis, we perceive and interact with our environment through the filters and categories of our language. The ability to talk about situations or characteristics of objects alters the way we perceive them \parencite{blasi2022over,kramsch2014language, whorfLanguage2012}.
In the field of cognitive science, the way in which culture including economy \parencite{fine2016world,henrich2001search}, values \parencite{inglehart2000world,triandis2018individualism} and language influence cognition is an object of intense investigation \parencite{barrett2020towards, henrichWeirdest2010}. 
The domain of color terms is a classic case-study for investigating culture-perception interaction because it has a directly measurable physiological correspondence.
It has been empirically demonstrated that when presented with the same color patches, people speaking different languages use different naming patterns. 

In a seminal work, Berlin and Kay subdivided color space into 330 Munsell Chips (Fig. \ref{fig:mode_prop_maps} upper panel) and ask participants to name these colors using "basic" color terms \parencite{berlinBasic1969}.
In the World Color Survey (WCS) \parencite{kayWorld2009} they examined 110 unwritten languages spoken solely of small scale societies.
From the investigation of these data they derived two main findings: First, different languages use different naming patterns - especially concerning the number of color terms. They claimed to have found a set of maximally eleven basic color categories. Languages make use of these terms or a subset of them down to three color terms. 
Second, languages using the same number of color terms exhibit similarities in the position of these categories in color space. This suggests universal rules shaping color terminology across languages.

\textcite{berlinBasic1969} suggested that exposure to global cultural influences has a leveling effect on the linguistic differences towards the eleven Basic Color Terms (BCT) in English vocabulary.
This is also why there was no large-scale, standardized, contemporary color-naming data collection at a global scale for industrialized languages. 
For some languages data sets exist but these investigations mostly concentrate on single languages and therefore differ in methodology and availability (for example \cite{lindseyColor2014,al-rasheedFurther2014, daviesBasic1994, kurikiModern2017, ozgenTurkish1998, thierryUnconscious2009, zollingerWhy1984,winawerRussian2007,berlinBasic1969}).

Globalization is a process having a heavy impact on informational flux and cultural exchange around the globe \parencite{barrett2020towards, pieterse2019globalization}. The work of this thesis is dedicated to the investigation of globalization effects on perception, particularly in the area of color terms.

With the recent publication of Large Language Models (LLM) e.g. GPT4 \parencite{achiam2023gpt} trained on all texts available online, another question can be addressed. If color perception is influenced by language and therefore conserved within text, are LLMs capable of reproducing language specific color naming patterns?

LLMs such as GPT4 \parencite{achiam2023gpt}, which are trained on extensive text corpora, have opened up a new line of inquiry.
%With the increase use of LLMs such as GPT4 \parencite{achiam2023gpt}, which are trained on extensive corpora of text accessible online, a novel line of inquiry emerges.
Since color perception is modulated by linguistic frameworks, and therefore embedded within textual corpora, can LLMs replicate language-specific color naming patterns?
In addition, the utilization of GPT4, a multilingual model subjected to extensive global cultural influences, may be considered as a proxy for a control group in assessing the impacts of globalization.

%In order to do justice to these impulses, 
To test this we collected data from human participants being native speakers of 25 distinct languages. Participants were recruited via the platform Lucid \parencite{lucid2024} , which served as a recruitment channel offering a non-traditional and diverse participant pool compared to traditional providers such as Amazon Mechanical Turk and Prolific.
In parallel, an analog dataset was created using the same procedure but engaging multilingual state-of-the-art LLMs (here GPT4; \cite{achiam2023gpt}) and Vision Language Models (VLMs, here GPT4V) as agents.

This thesis starts with an introduction describing the visual sense in general, ranging from light over the eye as the visual organ to the neurophysiological implementation in the brain stem and the visual cortex.
Then, I discuss the context information about the history and current state of knowledge in the field of color term research.